\newcommand{\matiere}{Physique : théorème de l'énergie cinétique}
% Préambule commun aux fiches d'exercices et de correction
\documentclass[11pt,a4paper]{article}
\usepackage[utf8]{inputenc}
\usepackage[T1]{fontenc}
\usepackage{lmodern}
\usepackage[french]{babel}
\usepackage[a4paper,margin=2.2cm,headheight=14pt]{geometry}
\usepackage{amsmath,amssymb}
\usepackage{icomma}
\usepackage[version=4]{mhchem}
\usepackage{chemfig}
\usepackage{tikz}
\usetikzlibrary{decorations.pathmorphing,arrows.meta,calc,patterns,angles,quotes}
\usepackage{enumitem}
\setlist[itemize]{label=\textbullet,nosep}
\setlist[enumerate]{nosep}
\usepackage{array}
\usepackage[most]{tcolorbox}
\usepackage{fancyhdr}
\usepackage{lastpage}
\usepackage{needspace}
\newcommand{\vv}[1]{\overrightarrow{#1}}

\setlength{\parindent}{0pt}
\setlength{\parskip}{5pt}
\setchemfig{atom sep=2.1em}

\pagestyle{fancy}
\fancyhf{}
\fancyhead[L]{\small Préparation du devoir}
\fancyhead[R]{\small\matiere}
\fancyfoot[C]{\small Page \thepage\ sur \pageref{LastPage}}
\renewcommand{\headrulewidth}{0.4pt}

% Titre d'un exercice
\newcommand{\exercice}[1]{%
  \par\bigskip\needspace{9\baselineskip}
  \begin{tcolorbox}[colback=gray!12,colframe=black,boxrule=0.5pt,arc=1pt,
    left=6pt,right=6pt,top=2pt,bottom=2pt]
  \textbf{#1}
  \end{tcolorbox}\par\smallskip}

% Titre d'une question dans la correction
\newcommand{\question}[1]{\par\medskip\textbf{#1}\par}

% Résultat encadré dans la correction
\newtcolorbox{resultat}{colback=gray!6,colframe=black,boxrule=0.4pt,arc=1pt,
  left=6pt,right=6pt,top=2pt,bottom=2pt}

% Données encadrées
\newtcolorbox{donnees}{colback=white,colframe=black!60,boxrule=0.4pt,arc=1pt,
  left=6pt,right=6pt,top=2pt,bottom=2pt}

% Page de titre
\newcommand{\entete}[2]{%
  \begin{center}
  {\LARGE\bfseries #1}\\[4pt]
  {\large #2}\\[2pt]
  {\normalsize Classe de Première}
  \end{center}}

\begin{document}

\entete{Série d'exercices}{Physique : théorème de l'énergie cinétique}

\vspace{2pt}
Nom et prénom : \dotfill \quad Classe : \dotfill

%-------------------------------------------------------------------------------
\exercice{Activité d'application}
\begin{minipage}[t]{0.56\textwidth}
\vspace{0pt}
Un objet de masse $m = 85$ kg descend une piste $ABC$ sans vitesse initiale avant de s'arrêter
en $C$. La partie $AB$ est un plan incliné d'un angle $\alpha$ par rapport à l'horizontale, sur
lequel les frottements sont négligeables. La partie $BC$ est horizontale.

\textbf{Données :} $\alpha = 30^\circ$ ; $AB = 20$ m ; $BC = 30$ m ; $g = 10$ N/kg.
\end{minipage}\hfill
\begin{minipage}[t]{0.40\textwidth}
\vspace{4pt}
\centering
\begin{tikzpicture}[scale=0.75]
  \coordinate (A) at (0,2.3); \coordinate (B) at (4,0); \coordinate (C) at (6.6,0);
  \coordinate (H) at (2,0);
  \draw[thick] (A) -- (B) -- (C);
  \draw[dashed] (B) -- (H);
  \pic[draw,"$\alpha$",angle radius=0.9cm,angle eccentricity=1.35] {angle = A--B--H};
  \fill (A) circle (1.5pt) node[above] {$A$};
  \fill (B) circle (1.5pt) node[below] {$B$};
  \fill (C) circle (1.5pt) node[below] {$C$};
\end{tikzpicture}
\end{minipage}

\begin{enumerate}
  \item Détermine la vitesse $V_B$ de l'objet au point $B$.
  \item Détermine la valeur $f$ des forces de frottement sur le parcours $BC$, supposées constantes.
\end{enumerate}

%-------------------------------------------------------------------------------
\exercice{Situation d'évaluation}
Un samedi matin des congés de Noël, tu effectues un voyage avec tes camarades de classe pour une
randonnée. La voiture chargée (voiture et passagers) a une masse totale $m = 1\,300$ kg.

Le conducteur démarre la voiture, aborde une côte avec la vitesse $v_1 = 3$ m.s$^{-1}$, puis atteint
son sommet avec la vitesse $v_2 = 12$ m.s$^{-1}$. La distance parcourue sur cette côte, dont la
ligne de plus grande pente fait un angle $\beta = 15^\circ$ avec le plan horizontal, est $L = 50$ m.

Du sommet de la côte, la voiture aborde une partie horizontale de la route en maintenant sa vitesse
constante sur une distance $d$, puis elle freine (moteur débrayé) et s'arrête sur une distance
$d' = 40$ m pour éviter de heurter un chien errant.

Durant tout le mouvement, les forces de frottement sont assimilées à une force unique $\vv{f}$ de
valeur constante $f = 780$ N. On prendra $g = 10$ N.kg$^{-1}$.

Tu es désigné par tes camarades pour montrer que les élèves de votre classe sont capables d'évaluer
les forces appliquées à la voiture.
\begin{enumerate}
  \item Énonce le théorème de l'énergie cinétique.
  \item Détermine la valeur $F$ de la force motrice $\vv{F}$ (force de propulsion exercée par le sol
        sur les roues de la voiture) :
  \begin{enumerate}[label=2.\arabic*.]
    \item durant son trajet sur la côte ;
    \item sur le plan horizontal, pendant que sa vitesse est constante.
  \end{enumerate}
  \item Détermine la valeur $F'$ de la force de freinage de la voiture.
\end{enumerate}

%-------------------------------------------------------------------------------
\clearpage
\exercice{Exercice 1}
\begin{minipage}[t]{0.58\textwidth}
\vspace{0pt}
Un petit objet $(S)$ de masse $m = 100$ g, assimilable à un point matériel, peut glisser à
l'intérieur d'une cuvette demi-sphérique de rayon $r = 0,5$ m. On le lâche sans vitesse initiale
du bord $A$ de cette cuvette. Il atteint le fond $B$ avec une vitesse $V_B$.
On prendra $g = 10$ N/kg.
\end{minipage}\hfill
\begin{minipage}[t]{0.38\textwidth}
\vspace{2pt}
\centering
\begin{tikzpicture}[scale=0.8]
  \draw[thick] (-2,0) arc (180:360:2);
  \draw[dashed] (0,0) -- node[right] {$r$} (0,-2);
  \fill (0,0) circle (1.2pt) node[above] {$O$};
  \fill (-2,0) circle (1.5pt) node[left] {$A$};
  \fill (0,-2) circle (1.5pt) node[below] {$B$};
  \fill[black] (-1.93,-0.35) rectangle ++(0.22,0.22);
\end{tikzpicture}
\end{minipage}

\begin{enumerate}
  \item Calcule la vitesse $V_B$ de l'objet au point $B$, en supposant les frottements négligeables.
  \item En réalité, l'objet atteint le fond $B$ avec la vitesse $V'_B = 2,5$ m.s$^{-1}$.
        Précise si l'objet est soumis à des forces de frottement. Justifie ta réponse.
  \item Détermine :
  \begin{enumerate}[label=3.\arabic*.]
    \item le travail total $W(\vv{f})$ des forces de frottement au cours du mouvement de $A$ à $B$ ;
    \item l'intensité $f$ de ces forces, supposée constante.
  \end{enumerate}
\end{enumerate}

%-------------------------------------------------------------------------------
\exercice{Exercice 2}
Recopie les groupes de mots ci-dessous dans l'ordre, de manière à obtenir une phrase correcte en
rapport avec le théorème de l'énergie cinétique :

\begin{center}
/ d'un solide / entre ces deux instants. / somme algébrique / La variation de l'énergie cinétique /
entre deux instants / des travaux de / toutes les forces extérieures appliquées / est égale à la /
\end{center}

%-------------------------------------------------------------------------------
\exercice{Exercice 3}
Un volant en fonte est constitué d'un cylindre homogène de rayon $R = 50$ cm et de hauteur
$h = 1$ m. La masse volumique de la fonte est $\rho = 7\,600$ kg.m$^{-3}$.

\emph{Rappel :} le moment d'inertie d'un cylindre homogène de masse $m$ et de rayon $R$ par rapport
à son axe est $J_\Delta = \frac{1}{2} m R^2$.
\begin{enumerate}
  \item Détermine le moment d'inertie $J_\Delta$ du cylindre par rapport à son axe.
  \item Calcule son énergie cinétique $E_C$ lorsqu'il tourne à la vitesse
        $N_1 = 1\,400$ tr.min$^{-1}$.
  \item Sa vitesse diminue jusqu'à $N_2 = 1\,300$ tr.min$^{-1}$ en $\Delta t = 4$ s.
        Calcule la puissance moyenne restituée par la diminution de l'énergie cinétique du volant.
\end{enumerate}

%-------------------------------------------------------------------------------
\exercice{Exercice 4}
Au cours d'une séance de révision dans ta salle d'étude, ton voisin de classe te propose
d'appliquer vos connaissances sur l'énergie cinétique pour déterminer la hauteur maximale atteinte
par un projectile lancé à partir du jouet de ton petit frère. Ce jouet est constitué d'un canon à
ressort de longueur à vide $l_0$, capable de lancer un petit projectile de masse $m$ vers le haut.
Ton voisin place le canon verticalement, comprime le ressort jusqu'à la longueur $l$, puis le lâche.

\textbf{Données :}
\begin{itemize}
  \item $m = 0,030$ kg ; $l_0 = 10$ cm ; $l = 5$ cm ; $g = 10$ N/kg ;
  \item une force d'intensité $F = 10$ N provoque un raccourcissement $x = 0,5$ cm du ressort ;
  \item pendant la détente du ressort, le poids du projectile est négligeable devant la force
        exercée par le ressort ; les frottements sont négligeables.
\end{itemize}
\begin{enumerate}
  \item Énonce le théorème de l'énergie cinétique.
  \item
  \begin{enumerate}[label=2.\arabic*.]
    \item Donne l'expression de la constante de raideur $k$ du ressort en fonction de $F$ et de $x$.
    \item Calcule la valeur de $k$.
  \end{enumerate}
  \item Détermine :
  \begin{enumerate}[label=3.\arabic*.]
    \item la vitesse $v_1$ du projectile à la sortie du canon ;
    \item la hauteur maximale atteinte par le projectile au-dessus de la sortie du canon.
  \end{enumerate}
\end{enumerate}

%-------------------------------------------------------------------------------
\exercice{Exercice 5 (Première C)}
Un groupe d'élèves de Première scientifique se propose de déterminer la compression du ressort
utilisé dans un jeu d'enfant. La piste de jeu comporte une partie $AB$ en forme d'arc de cercle de
centre $O$ et de rayon $R$, prolongée par une partie horizontale $BCD$. Un ressort de raideur $k$,
dont l'extrémité libre est en $C$, est fixé sur la partie horizontale (voir la figure).

\begin{center}
\begin{tikzpicture}[scale=1.15]
  \coordinate (O) at (0,0); \coordinate (A) at (210:2); \coordinate (B) at (0,-2);
  \coordinate (C) at (3.6,-2); \coordinate (D) at (4.3,-2);
  \draw[thick] (A) arc (210:270:2) -- (7,-2);
  \draw[dashed] (O) -- (A); \draw[dashed] (O) -- (B);
  \pic[draw,"$\alpha$",angle radius=0.6cm,angle eccentricity=1.4] {angle = A--O--B};
  \fill (O) circle (1.3pt) node[above] {$O$};
  \fill (A) circle (1.5pt) node[below left] {$A$};
  \fill (B) circle (1.5pt) node[below] {$B$};
  \fill (240:2) circle (1.5pt) node[below left] {$M$};
  \fill (1.8,-2) circle (1.5pt) node[below] {$M'$};
  \node[below] at (C) {$C$}; \node[below] at (D) {$D$};
  \draw (D) -- ++(0,0.12);
  \draw[decorate,decoration={coil,aspect=0.5,segment length=2mm,amplitude=1.8mm}]
       (3.6,-1.8) -- (6.6,-1.8);
  \fill[pattern=north east lines] (6.6,-2) rectangle (6.85,-1.35);
  \draw (6.6,-2) -- (6.6,-1.35);
  \shade[ball color=black!70] ($(A)+(0,1.2)$) circle (0.12);
  \draw[<->] ($(A)+(-0.9,0)$) -- node[left] {$H$} ($(A)+(-0.9,1.2)$);
  \draw[densely dotted] (A) -- ++(-1.05,0);
  \draw[densely dotted] ($(A)+(0,1.2)$) -- ++(-1.05,0);
\end{tikzpicture}
\end{center}

Pour jouer, on laisse tomber sans vitesse initiale une bille ponctuelle de masse $m$, située à la
hauteur $H$ au-dessus du point $A$. La bille arrive en $A$ et se déplace alors sur la piste $ABC$.
Elle vient heurter l'extrémité libre du ressort, qui se comprime d'une longueur $x$, et elle
s'arrête en $D$. Au cours de son mouvement sur la piste, la bille est soumise à des forces de
frottement d'intensité $f$ constante, qui s'annulent au-delà du point $C$. On néglige la résistance
de l'air pendant la chute.

\textbf{Données :} $m = 30$ g ; $k = 20$ N.m$^{-1}$ ; $g = 10$ N/kg ; $H = 50$ cm ; $R = 50$ cm ;
$\alpha = 60^\circ$ ; $BC = 50$ cm ; $f = 4,3 \times 10^{-2}$ N.

Tu es désigné pour rédiger le compte rendu.
\begin{enumerate}
  \item Énonce le théorème de l'énergie cinétique.
  \item Calcule l'énergie cinétique de la bille lorsqu'elle arrive au point $A$.
  \item Représente les forces qui s'exercent sur la bille :
  \begin{enumerate}[label=3.\arabic*.]
    \item au point $M$ de l'arc $AB$ et au point $M'$ de la partie $BC$ ;
    \item lorsqu'elle est en contact avec le ressort comprimé de $x$.
  \end{enumerate}
  \item Calcule $x$, puis explique ce qui se passe après l'instant d'arrêt en $D$.
\end{enumerate}

\end{document}
